% ====================================================================
% ФАЙЛ: tasks/06_integrated/kim_task18_all_variants.tex
% ТЕМА: ИНТЕГРИРОВАННЫЕ ЗАДАНИЯ №18 (30 ВАРИАНТОВ ЕГЭ 2026)
% ====================================================================

% === ВАРИАНТ 1 ===
% Тег: #МЕХАНИКА #МКТ #ЭЛЕКТРОДИНАМИКА #КВАНТОВАЯ_ФИЗИКА #КИМ_18 #ВАРИАНТ_01
\begin{taskbasic}[code={\label{task:kim18_v01}}]
Выберите все верные утверждения о физических явлениях, величинах и закономерностях. Запишите цифры, под которыми они указаны.

1) Работа силы тяжести по перемещению тела между двумя заданными точками не зависит от длины соединяющей их траектории.\\
2) При прочих равных условиях диффузия протекает в жидкостях значительно медленнее, чем в твёрдых телах.\\
3) Весь электростатический заряд проводника сосредоточен на его поверхности.\\
4) Свободные электромагнитные колебания являются гармоническими, если электрический заряд на обкладках конденсатора с течением времени меняется по закону синуса или косинуса.\\
5) Атомы изотопов одного и того же химического элемента различаются числом протонов.

\kim{18}
\end{taskbasic}
\answer{134}

% === ВАРИАНТ 2 ===
% Тег: #МЕХАНИКА #МКТ #ЭЛЕКТРОДИНАМИКА #КВАНТОВАЯ_ФИЗИКА #КИМ_18 #ВАРИАНТ_02
\begin{taskbasic}[code={\label{task:kim18_v02}}]
Выберите все верные утверждения о физических явлениях, величинах и закономерностях. Запишите цифры, под которыми они указаны.

1) Работа силы тяжести по перемещению тела между двумя заданными точками зависит от длины соединяющей их траектории.\\
2) При прочих равных условиях диффузия протекает в твёрдых телах значительно медленнее, чем в жидкостях.\\
3) Весь электростатический заряд проводника сосредоточен в его центре масс.\\
4) Свободные электромагнитные колебания являются затухающими, если электрический заряд на обкладках конденсатора с течением времени меняется по закону синуса или косинуса.\\
5) Атомы изотопов одного и того же химического элемента различаются числом нейтронов.

\kim{18}
\end{taskbasic}
\answer{25}

% === ВАРИАНТ 3 ===
% Тег: #МЕХАНИКА #МКТ #ЭЛЕКТРОДИНАМИКА #КВАНТОВАЯ_ФИЗИКА #КИМ_18 #ВАРИАНТ_03
\begin{taskbasic}[code={\label{task:kim18_v03}}]
Выберите все верные утверждения о физических явлениях, величинах и закономерностях. Запишите цифры, под которыми они указаны.

1) Модуль максимальной силы трения покоя обратно пропорционален силе нормального давления на опору.\\
2) В ходе процесса кристаллизации тела его температура не меняется, а внутренняя энергия системы «жидкость + твёрдое тело» уменьшается.\\
3) При протекании электрического тока в любых средах обязательно наблюдается нагревание проводника.\\
4) Индукционный ток своим магнитным полем всегда противодействует тому изменению магнитного потока, которым он вызван.\\
5) При электронном $\beta$-распаде зарядовое число ядра остаётся практически неизменным.

\kim{18}
\end{taskbasic}
\answer{24}

% === ВАРИАНТ 4 ===
% Тег: #МЕХАНИКА #МКТ #ЭЛЕКТРОДИНАМИКА #КВАНТОВАЯ_ФИЗИКА #КИМ_18 #ВАРИАНТ_04
\begin{taskbasic}[code={\label{task:kim18_v04}}]
Выберите все верные утверждения о физических явлениях, величинах и закономерностях. Запишите цифры, под которыми они указаны.

1) Модуль максимальной силы трения покоя прямо пропорционален силе нормального давления на опору.\\
2) В ходе процесса кристаллизации тела его температура не меняется, а внутренняя энергия системы «жидкость + твёрдое тело» возрастает.\\
3) При протекании электрического тока в проводнике вокруг него возникает магнитное поле.\\
4) Индукционный ток своим магнитным полем всегда увеличивает магнитный поток, изменением которого он вызван.\\
5) При электронном $\beta$-распаде массовое число ядра остаётся практически неизменным.

\kim{18}
\end{taskbasic}
\answer{135}

% === ВАРИАНТ 5 ===
% Тег: #МЕХАНИКА #МКТ #ЭЛЕКТРОДИНАМИКА #КВАНТОВАЯ_ФИЗИКА #КИМ_18 #ВАРИАНТ_05
\begin{taskbasic}[code={\label{task:kim18_v05}}]
Выберите все верные утверждения о физических явлениях, величинах и закономерностях. Запишите цифры, под которыми они указаны.

1) Сила Архимеда, действующая на полностью погружённое в жидкость тело, прямо пропорциональна плотности жидкости.\\
2) Теплопередача путём теплопроводности происходит за счёт переноса вещества струями и потоками.\\
3) Работа электростатического поля по перемещению заряда между двумя заданными точками не зависит от формы траектории.\\
4) Индукционный ток возникает в разомкнутом проводящем контуре при изменении магнитного потока, пронизывающего контур.\\
5) В планетарной модели атома число нейтронов в ядре равно числу электронов в электронной оболочке нейтрального атома.

\kim{18}
\end{taskbasic}
\answer{13}

% === ВАРИАНТ 6 ===
% Тег: #МЕХАНИКА #МКТ #ЭЛЕКТРОДИНАМИКА #КВАНТОВАЯ_ФИЗИКА #КИМ_18 #ВАРИАНТ_06
\begin{taskbasic}[code={\label{task:kim18_v06}}]
Выберите все верные утверждения о физических явлениях, величинах и закономерностях. Запишите цифры, под которыми они указаны.

1) Сила Архимеда, действующая на полностью погружённое в жидкость тело, прямо пропорциональна его плотности.\\
2) Теплопередача путём конвекции происходит за счёт переноса вещества струями и потоками.\\
3) Сила Лоренца, действующая на заряженную частицу, параллельна её скорости.\\
4) Индукционный ток возникает в замкнутом проводящем контуре при изменении магнитного потока, пронизывающего контур.\\
5) В планетарной модели атома число протонов в ядре равно числу электронов в электронной оболочке нейтрального атома.

\kim{18}
\end{taskbasic}
\answer{245}

% === ВАРИАНТ 7 ===
% Тег: #МЕХАНИКА #МКТ #ЭЛЕКТРОДИНАМИКА #ОПТИКА #КВАНТОВАЯ_ФИЗИКА #КИМ_18 #ВАРИАНТ_07
\begin{taskbasic}[code={\label{task:kim18_v07}}]
Выберите все верные утверждения о физических явлениях, величинах и закономерностях. Запишите цифры, под которыми они указаны.

1) Если скорость тела с течением времени увеличивается, то вектор ускорения направлен противоположно вектору скорости тела.\\
2) В ходе процесса кристаллизации жидкости внутренняя энергия системы «жидкость + твёрдое тело» увеличивается, а температура этой системы остаётся неизменной.\\
3) Энергия электрического поля конденсатора, заряженного и отключённого от источника, не меняется при изменении электроёмкости конденсатора.\\
4) При переходе электромагнитной волны из оптически менее плотной в оптически более плотную среду длина волны уменьшается.\\
5) В планетарной модели атома число протонов в ядре равно числу электронов в электронной оболочке нейтрального атома.

\kim{18}
\end{taskbasic}
\answer{45}

% === ВАРИАНТ 8 ===
% Тег: #МЕХАНИКА #МКТ #ЭЛЕКТРОДИНАМИКА #ОПТИКА #КВАНТОВАЯ_ФИЗИКА #КИМ_18 #ВАРИАНТ_08
\begin{taskbasic}[code={\label{task:kim18_v08}}]
Выберите все верные утверждения о физических явлениях, величинах и закономерностях. Запишите цифры, под которыми они указаны.

1) Если скорость тела, движущегося по прямой, с течением времени увеличивается, то вектор ускорения сонаправлен с вектором скорости тела.\\
2) В ходе процесса кристаллизации жидкости внутренняя энергия системы «жидкость + твёрдое тело» уменьшается, а температура этой системы остаётся неизменной.\\
3) Энергия электрического поля конденсатора, заряженного и отключённого от источника, уменьшается при увеличении электроёмкости конденсатора.\\
4) При переходе электромагнитной волны из оптически менее плотной в оптически более плотную среду длина волны увеличивается.\\
5) В планетарной модели атома число нейтронов в ядре равно числу электронов в электронной оболочке нейтрального атома.

\kim{18}
\end{taskbasic}
\answer{123}

% === ВАРИАНТ 9 ===
% Тег: #МЕХАНИКА #МКТ #ЭЛЕКТРОДИНАМИКА #КВАНТОВАЯ_ФИЗИКА #КИМ_18 #ВАРИАНТ_09
\begin{taskbasic}[code={\label{task:kim18_v09}}]
Выберите все верные утверждения о физических явлениях, величинах и закономерностях. Запишите цифры, под которыми они указаны.

1) При равномерном движении по окружности вектор скорости материальной точки не изменяется.\\
2) При изотермическом сжатии постоянной массы одноатомного идеального газа его внутренняя энергия остаётся неизменной.\\
3) В электростатическом поле работа по перемещению заряда между двумя точками поля не зависит от формы траектории, а определяется только начальным и конечным положениями заряда.\\
4) Силой Лоренца называют силу, с которой магнитное поле действует на движущуюся в нём заряженную частицу.\\
5) При электронном $\beta$-распаде заряд ядра уменьшается на величину одного элементарного заряда.

\kim{18}
\end{taskbasic}
\answer{234}

% === ВАРИАНТ 10 ===
% Тег: #МЕХАНИКА #МКТ #ЭЛЕКТРОДИНАМИКА #КВАНТОВАЯ_ФИЗИКА #КИМ_18 #ВАРИАНТ_10
\begin{taskbasic}[code={\label{task:kim18_v10}}]
Выберите все верные утверждения о физических явлениях, величинах и закономерностях. Запишите цифры, под которыми они указаны.

1) При прямолинейном равноускоренном движении тела модуль его ускорения со временем не изменяется.\\
2) При изохорном нагревании газа постоянной массы его давление уменьшается.\\
3) Поверхность проводника любой формы в электростатическом поле является эквипотенциальной.\\
4) Сила Лоренца всегда перпендикулярна вектору скорости заряженной частицы, движущейся в однородном магнитном поле.\\
5) При $\alpha$-распаде массовое число ядра-продукта увеличивается на 4 по сравнению с исходным ядром.

\kim{18}
\end{taskbasic}
\answer{134}

% ====================================================================
% ФАЙЛ: tasks/06_integrated/kim_task18_part2.tex
% ТЕМА: ИНТЕГРИРОВАННЫЕ ЗАДАНИЯ №18 (ДОПОЛНИТЕЛЬНАЯ ВЫРЕЗКА)
% ====================================================================

% === ЗАДАЧА 1 ===
% Тег: #МЕХАНИКА #МКТ #ЭЛЕКТРОДИНАМИКА #ОПТИКА #КВАНТОВАЯ_ФИЗИКА #КИМ_18
\begin{taskbasic}[code={\label{task:kim18_img1}}]
Выберите все верные утверждения о физических явлениях, величинах и закономерностях. Запишите цифры, под которыми они указаны.

1) При равномерном прямолинейном движении за любые равные промежутки времени тело совершает одинаковые перемещения.\\
2) Теплопередача путём излучения происходит за счёт переноса вещества струями и потоками.\\
3) Одноимённые точечные электрические заряды отталкиваются друг от друга.\\
4) Дифракция волн особенно хорошо наблюдается в тех случаях, когда размеры препятствий существенно превышают длину волны.\\
5) Изотопы химического элемента имеют в ядре одинаковое число протонов и разное число нейтронов.

\kim{18}
\end{taskbasic}
\answer{135}

% === ЗАДАЧА 2 ===
% Тег: #МЕХАНИКА #МКТ #ЭЛЕКТРОДИНАМИКА #ОПТИКА #КВАНТОВАЯ_ФИЗИКА #КИМ_18
\begin{taskbasic}[code={\label{task:kim18_img2}}]
Выберите все верные утверждения о физических явлениях, величинах и закономерностях. Запишите цифры, под которыми они указаны.

1) При равноускоренном прямолинейном движении за любые равные промежутки времени тело совершает одинаковые перемещения.\\
2) Теплопередача путём конвекции происходит за счёт переноса вещества струями и потоками.\\
3) Одноимённые точечные электрические заряды притягиваются друг к другу.\\
4) Дифракция волн хорошо наблюдается в тех случаях, когда размеры препятствий сравнимы с длиной волны.\\
5) Изотопы химического элемента имеют в ядре одинаковое число нейтронов и разное число протонов.

\kim{18}
\end{taskbasic}
\answer{24}

% === ЗАДАЧА 3 ===
% Тег: #МЕХАНИКА #МКТ #ЭЛЕКТРОДИНАМИКА #ОПТИКА #КВАНТОВАЯ_ФИЗИКА #КИМ_18
\begin{taskbasic}[code={\label{task:kim18_img3}}]
Выберите все верные утверждения о физических явлениях, величинах и закономерностях. Запишите цифры, под которыми они указаны.

1) Модуль сил гравитационного взаимодействия двух материальных точек обратно пропорционален квадрату расстояния между ними.\\
2) Давление насыщенного пара увеличивается с ростом абсолютной температуры пара и не зависит от его объёма.\\
3) В однородном электростатическом поле работа силы, действующей на заряд со стороны поля, при перемещении заряда между двумя заданными точками прямо пропорциональна длине выбранной траектории.\\
4) При переходе электромагнитной волны из оптически менее плотной в оптически более плотную среду частота волны остаётся неизменной.\\
5) При распространении света проявляются только его корпускулярные свойства, а при взаимодействии с веществом --- только волновые.

\kim{18}
\end{taskbasic}
\answer{124}

% === ЗАДАЧА 4 ===
% Тег: #МЕХАНИКА #МКТ #ЭЛЕКТРОДИНАМИКА #ОПТИКА #КВАНТОВАЯ_ФИЗИКА #КИМ_18
\begin{taskbasic}[code={\label{task:kim18_img4}}]
Выберите все верные утверждения о физических явлениях, величинах и закономерностях. Запишите цифры, под которыми они указаны.

1) Модуль сил гравитационного взаимодействия двух материальных точек прямо пропорционален квадрату расстояния между ними.\\
2) Давление насыщенного пара увеличивается с ростом абсолютной температуры пара и зависит от его объёма.\\
3) В электростатическом поле работа силы, действующей на заряд со стороны поля, при перемещении заряда между двумя заданными точками не зависит от траектории.\\
4) При переходе электромагнитной волны из оптически менее плотной в оптически более плотную среду длина волны остаётся неизменной.\\
5) При распространении и взаимодействии света с веществом проявляются как волновые, так и его корпускулярные свойства.

\kim{18}
\end{taskbasic}
\answer{35}

% === ЗАДАЧА 5 ===
% Тег: #МЕХАНИКА #МКТ #ЭЛЕКТРОДИНАМИКА #КВАНТОВАЯ_ФИЗИКА #КИМ_18
\begin{taskbasic}[code={\label{task:kim18_img5}}]
Выберите все верные утверждения о физических явлениях, величинах и закономерностях. Запишите цифры, под которыми они указаны.

1) Импульсом силы называется величина, равная произведению массы тела на его ускорение.\\
2) В изотермическом процессе для постоянной массы газа отношение объёма газа к его давлению остаётся постоянным.\\
3) Модуль сил взаимодействия двух неподвижных точечных заряженных тел обратно пропорционален квадрату расстояния между ними.\\
4) Период свободных электромагнитных колебаний в идеальном колебательном контуре увеличивается прямо пропорционально увеличению электроёмкости конденсатора.\\
5) В планетарной модели атома число протонов в ядре равно числу электронов в электронной оболочке нейтрального атома.

\kim{18}
\end{taskbasic}
\answer{35}

% === ЗАДАЧА 6 ===
% Тег: #МЕХАНИКА #МКТ #ЭЛЕКТРОДИНАМИКА #КВАНТОВАЯ_ФИЗИКА #КИМ_18
\begin{taskbasic}[code={\label{task:kim18_img6}}]
Выберите все верные утверждения о физических явлениях, величинах и закономерностях. Запишите цифры, под которыми они указаны.

1) Импульсом силы называется величина, равная произведению силы, действующей на тело, на время её действия.\\
2) В изотермическом процессе для постоянной массы газа произведение объёма газа на его давление остаётся постоянным.\\
3) Модуль сил взаимодействия двух неподвижных точечных заряженных тел обратно пропорционален расстоянию между ними.\\
4) Частота свободных электромагнитных колебаний в идеальном колебательном контуре увеличивается с уменьшением индуктивности катушки.\\
5) В планетарной модели атома число нейтронов в ядре равно числу электронов в электронной оболочке нейтрального атома.

\kim{18}
\end{taskbasic}
\answer{124}

% === ЗАДАЧА 7 ===
% Тег: #МЕХАНИКА #МКТ #ЭЛЕКТРОДИНАМИКА #КВАНТОВАЯ_ФИЗИКА #КИМ_18
\begin{taskbasic}[code={\label{task:kim18_img7}}]
Выберите все верные утверждения о физических явлениях, величинах и закономерностях. Запишите цифры, под которыми они указаны.

1) Период гармонических колебаний колебательной системы прямо пропорционален частоте её колебаний.\\
2) Внутренняя энергия постоянной массы одноатомного идеального газа уменьшается при понижении абсолютной температуры газа.\\
3) Изначально незаряженные тела в процессе электризации трением приобретают равные по модулю и противоположные по знаку заряды.\\
4) При монотонном изменении магнитного потока, пронизывающего площадку, ограниченную разомкнутым проводящим контуром, в контуре возникает индукционный ток.\\
5) В планетарной модели атома число нейтронов в ядре равно числу электронов в электронной оболочке нейтрального атома.

\kim{18}
\end{taskbasic}
\answer{23}

% === ЗАДАЧА 8 ===
% Тег: #МЕХАНИКА #МКТ #ЭЛЕКТРОДИНАМИКА #КВАНТОВАЯ_ФИЗИКА #КИМ_18
\begin{taskbasic}[code={\label{task:kim18_img8}}]
Выберите все верные утверждения о физических явлениях, величинах и закономерностях. Запишите цифры, под которыми они указаны.

1) Период гармонических колебаний колебательной системы обратно пропорционален частоте её колебаний.\\
2) Внутренняя энергия постоянной массы одноатомного идеального газа увеличивается при понижении абсолютной температуры газа.\\
3) Изначально незаряженные тела в процессе электризации трением приобретают равные по модулю и одинаковые по знаку заряды.\\
4) В замкнутом проводящем контуре при изменении магнитного потока, пронизывающего контур, возникает индукционный ток.\\
5) В планетарной модели атома число протонов в ядре равно числу электронов в электронной оболочке нейтрального атома.

\kim{18}
\end{taskbasic}
\answer{145}

% === ЗАДАЧА 9 ===
% Тег: #МЕХАНИКА #МКТ #ЭЛЕКТРОДИНАМИКА #КВАНТОВАЯ_ФИЗИКА #КИМ_18
\begin{taskbasic}[code={\label{task:kim18_img9}}]
Выберите все верные утверждения о физических явлениях, величинах и закономерностях. Запишите цифры, под которыми они указаны.

1) Модуль сил гравитационного взаимодействия двух тел прямо пропорционален квадрату расстояния между этими телами.\\
2) Теплопередача путём конвекции происходит за счёт переноса энергии струями и потоками жидкости или газа.\\
3) Модуль сил взаимодействия двух неподвижных точечных заряженных тел не зависит от свойств среды между ними.\\
4) Период свободных колебаний в идеальном колебательном контуре увеличивается прямо пропорционально увеличению индуктивности катушки.\\
5) При $\alpha$-распаде масса ядра уменьшается примерно на четыре атомных единицы массы.

\kim{18}
\end{taskbasic}
\answer{25}

% === ЗАДАЧА 10 ===
% Тег: #МЕХАНИКА #МКТ #ЭЛЕКТРОДИНАМИКА #КВАНТОВАЯ_ФИЗИКА #КИМ_18
\begin{taskbasic}[code={\label{task:kim18_img10}}]
Выберите все верные утверждения о физических явлениях, величинах и закономерностях. Запишите цифры, под которыми они указаны.

1) Модуль сил гравитационного взаимодействия двух тел обратно пропорционален квадрату расстояния между этими телами.\\
2) Теплопередача путём теплопроводности происходит за счёт переноса энергии струями и потоками жидкости или газа.\\
3) Модуль сил взаимодействия двух неподвижных точечных заряженных тел зависит от свойств среды между ними.\\
4) Частота свободных колебаний в идеальном колебательном контуре уменьшается с увеличением индуктивности катушки.\\
5) При $\alpha$-распаде масса ядра уменьшается примерно на две атомные единицы массы.

\kim{18}
\end{taskbasic}
\answer{134}

% ====================================================================
% ФАЙЛ: tasks/06_integrated/kim_task18_part3.tex
% ТЕМА: ИНТЕГРИРОВАННЫЕ ЗАДАНИЯ №18 (ЗАВЕРШАЮЩИЙ ПАКЕТ)
% ====================================================================

% === ЗАДАЧА 1 ===
% Тег: #МЕХАНИКА #МКТ #ЭЛЕКТРОДИНАМИКА #ОПТИКА #КВАНТОВАЯ_ФИЗИКА #КИМ_18
\begin{taskbasic}[code={\label{task:kim18_img11}}]
Выберите все верные утверждения о физических явлениях, величинах и закономерностях. Запишите цифры, под которыми они указаны.

1) При вынужденных механических колебаниях в колебательной системе резонанс возникает в том случае, если собственная частота колебаний системы совпадает с частотой изменения внешней силы.\\
2) В процессе изохорного нагревания постоянной массы газа давление газа уменьшается.\\
3) Поверхность проводника, находящегося в электростатическом поле, является эквипотенциальной.\\
4) При преломлении света при переходе из одной среды в другую изменяются скорость волны и длина волны, а её частота остаётся неизменной.\\
5) Энергия связи ядра равна той энергии, которая выделяется при реакции синтеза ядра из ядер более лёгких изотопов.

\kim{18}
\end{taskbasic}
\answer{134}

% === ЗАДАЧА 2 ===
% Тег: #МЕХАНИКА #МКТ #ЭЛЕКТРОДИНАМИКА #ОПТИКА #КВАНТОВАЯ_ФИЗИКА #КИМ_18
\begin{taskbasic}[code={\label{task:kim18_img12}}]
Выберите все верные утверждения о физических явлениях, величинах и закономерностях. Запишите цифры, под которыми они указаны.

1) При вынужденных механических колебаниях в колебательной системе резонанс возникает в том случае, если собственная частота колебаний системы превышает частоту изменения внешней силы.\\
2) В процессе изохорного нагревания постоянной массы газа давление газа увеличивается.\\
3) Поверхность проводника, находящегося в электростатическом поле, не является эквипотенциальной.\\
4) При преломлении света при переходе из одной среды в другую изменяются скорость волны и частота, а длина её волны остаётся неизменной.\\
5) Энергия связи ядра равна той энергии, которую необходимо затратить для того, чтобы разделить это ядро на отдельные протоны и нейтроны.

\kim{18}
\end{taskbasic}
\answer{25}

% === ЗАДАЧА 3 ===
% Тег: #МЕХАНИКА #МКТ #ЭЛЕКТРОДИНАМИКА #КВАНТОВАЯ_ФИЗИКА #КИМ_18
\begin{taskbasic}[code={\label{task:kim18_img13}}]
Выберите все верные утверждения о физических явлениях, величинах и закономерностях. Запишите цифры, под которыми они указаны.

1) Громкость звука определяется частотой колебаний.\\
2) Температура плавления кристаллических тел зависит от их массы.\\
3) В цепи постоянного тока на всех параллельно соединённых резисторах напряжение одинаково.\\
4) Скорость распространения радиоволн в вакууме равна скорости света в вакууме.\\
5) При $\beta$-распаде ядра образуется ион нового элемента и ядро атома гелия.

\kim{18}
\end{taskbasic}
\answer{34}

% === ЗАДАЧА 4 ===
% Тег: #МЕХАНИКА #МКТ #ЭЛЕКТРОДИНАМИКА #КВАНТОВАЯ_ФИЗИКА #КИМ_18
\begin{taskbasic}[code={\label{task:kim18_img14}}]
Выберите все верные утверждения о физических явлениях, величинах и закономерностях. Запишите цифры, под которыми они указаны.

1) Громкость звука определяется амплитудой звуковых колебаний.\\
2) Температура плавления кристаллических тел не зависит от их массы.\\
3) В цепи постоянного тока на всех последовательно соединённых резисторах напряжение одинаково.\\
4) Скорость распространения радиоволн в воде больше скорости света в вакууме.\\
5) При $\alpha$-распаде ядра образуется ядро нового элемента и ядро атома гелия.

\kim{18}
\end{taskbasic}
\answer{125}

% === ЗАДАЧА 5 ===
% Тег: #МЕХАНИКА #МКТ #ЭЛЕКТРОДИНАМИКА #ОПТИКА #КВАНТОВАЯ_ФИЗИКА #КИМ_18
\begin{taskbasic}[code={\label{task:kim18_img15}}]
Выберите все верные утверждения о физических явлениях, величинах и закономерностях. Запишите цифры, под которыми они указаны.

1) Давление столба жидкости на дно сосуда обратно пропорционально её плотности.\\
2) Удельная теплота плавления вещества показывает, какое количество теплоты необходимо сообщить 1 кг вещества, находящемуся при любой температуре, чтобы его расплавить.\\
3) В процессе электризации трением два первоначально незаряженных тела приобретают разноимённые и одинаковые по модулю заряды.\\
4) При переходе света из оптически более плотной среды в оптически менее плотную среду угол падения меньше угла преломления.\\
5) При $\alpha$-распаде ядра выполняется закон сохранения электрического заряда, но не выполняется закон сохранения импульса.

\kim{18}
\end{taskbasic}
\answer{34}

% === ЗАДАЧА 6 ===
% Тег: #МЕХАНИКА #МКТ #ЭЛЕКТРОДИНАМИКА #ОПТИКА #КВАНТОВАЯ_ФИЗИКА #КИМ_18
\begin{taskbasic}[code={\label{task:kim18_img16}}]
Выберите все верные утверждения о физических явлениях, величинах и закономерностях. Запишите цифры, под которыми они указаны.

1) Давление столба жидкости на дно сосуда прямо пропорционально её плотности.\\
2) Удельная теплота плавления вещества показывает, какое количество теплоты необходимо сообщить 1 кг вещества, находящемуся при температуре плавления, чтобы его расплавить.\\
3) В процессе электризации трением два первоначально незаряженных тела приобретают разноимённые и различные по модулю заряды.\\
4) При переходе света из оптически более плотной среды в оптически менее плотную среду угол падения больше угла преломления.\\
5) При $\alpha$-распаде ядра выполняются закон сохранения электрического заряда и закон сохранения импульса.

\kim{18}
\end{taskbasic}
\answer{125}

% === ЗАДАЧА 7 ===
% Тег: #МЕХАНИКА #МКТ #ЭЛЕКТРОДИНАМИКА #ОПТИКА #КВАНТОВАЯ_ФИЗИКА #КИМ_18
\begin{taskbasic}[code={\label{task:kim18_img17}}]
Выберите все верные утверждения о физических явлениях, величинах и закономерностях. Запишите цифры, под которыми они указаны.

1) При прохождении математическим маятником положения равновесия центростремительное ускорение его груза максимально.\\
2) Удельная теплоёмкость вещества показывает, какое количество теплоты необходимо сообщить 1 кг вещества для его плавления.\\
3) При помещении проводника в электростатическое поле наблюдается явление электростатической индукции.\\
4) При преломлении света, падающего из среды с меньшим показателем преломления в среду с большим показателем преломления, угол падения меньше угла преломления.\\
5) При $\beta$-распаде ядра выполняются законы сохранения энергии и электрического заряда, но не выполняется закон сохранения импульса.

\kim{18}
\end{taskbasic}
\answer{13}

% === ЗАДАЧА 8 ===
% Тег: #МЕХАНИКА #МКТ #ЭЛЕКТРОДИНАМИКА #ОПТИКА #КВАНТОВАЯ_ФИЗИКА #КИМ_18
\begin{taskbasic}[code={\label{task:kim18_img18}}]
Выберите все верные утверждения о физических явлениях, величинах и закономерностях. Запишите цифры, под которыми они указаны.

1) При прохождении математическим маятником положения равновесия центростремительное ускорение его груза равно нулю.\\
2) Удельная теплоёмкость вещества показывает, какое количество теплоты необходимо сообщить 1 кг вещества для его нагревания на 1 К.\\
3) При помещении проводника в электростатическое поле наблюдается явление электромагнитной индукции.\\
4) При преломлении света, падающего из среды с меньшим показателем преломления в среду с большим показателем преломления, угол падения больше угла преломления.\\
5) При $\beta$-распаде ядра выполняются законы сохранения энергии и электрического заряда и закон сохранения импульса.

\kim{18}
\end{taskbasic}
\answer{245}

% === ЗАДАЧА 9 ===
% Тег: #МЕХАНИКА #МКТ #ЭЛЕКТРОДИНАМИКА #ОПТИКА #КВАНТОВАЯ_ФИЗИКА #КИМ_18
\begin{taskbasic}[code={\label{task:kim18_img19}}]
Выберите все верные утверждения о физических явлениях, величинах и закономерностях. Запишите цифры, под которыми они указаны.

1) При равномерном прямолинейном движении за любые равные промежутки времени тело совершает одинаковые перемещения.\\
2) Средняя кинетическая энергия поступательного теплового движения молекул газа обратно пропорциональна абсолютной температуре газа.\\
3) В однородном электростатическом поле работа по перемещению заряда между двумя точками не зависит от траектории.\\
4) При переходе электромагнитной волны из оптически менее плотной в оптически более плотную среду частота волны уменьшается.\\
5) При электронном $\beta$-распаде массовое число ядра остаётся неизменным.

\kim{18}
\end{taskbasic}
\answer{135}

% === ЗАДАЧА 10 ===
% Тег: #МЕХАНИКА #МКТ #ЭЛЕКТРОДИНАМИКА #ОПТИКА #КВАНТОВАЯ_ФИЗИКА #КИМ_18
\begin{taskbasic}[code={\label{task:kim18_img20}}]
Выберите все верные утверждения о физических явлениях, величинах и закономерностях. Запишите цифры, под которыми они указаны.

1) При равномерном движении по окружности за любые равные промежутки времени тело совершает одинаковые перемещения.\\
2) Средняя кинетическая энергия поступательного теплового движения молекул газа прямо пропорциональна абсолютной температуре газа.\\
3) В неоднородном электростатическом поле работа по перемещению заряда между двумя точками зависит от траектории.\\
4) При переходе электромагнитной волны из оптически менее плотной в оптически более плотную среду частота волны не изменяется.\\
5) При электронном $\beta$-распаде массовое число ядра уменьшается.

\kim{18}
\end{taskbasic}
\answer{24}